\documentclass[11pt]{amsart}

\usepackage{amsmath, amssymb, amsthm}
\usepackage{mathtools}
\usepackage{thmtools}
\usepackage{enumitem}
\usepackage{tikz-cd}
\usepackage{hyperref}
\usepackage[nameinlink]{cleveref}

% % Theorem environments
\declaretheorem[numberwithin=section]{theorem}
\declaretheorem[style=plain, sibling=theorem]{lemma, proposition, corollary}
\declaretheorem[style=definition]{definition, example, condition}
\declaretheorem[style=remark]{remark, note, observation, todo}


\title{April 18, 2026}

\begin{document}
\maketitle

I reviewed the feedback file for the daily writing.
The arXiv check file is also generated, but I did not check it since there were lots of duplicated files.
Instead of checking, I updated the sourcode managing arXiv files to deal with duplicated papers.

\par\vspace{1em}

I felt that it was a time to put things I learned into together.
The time I spent will teach me where am I going.
Specifically, the relation between hyperbolic geometry, $\mathrm{SL}_2$, its coordinatering, and the quantum group $\mathrm{SL}_q(2)$ as its quantization, and the volume conjecture.
The first thing to summarize is that the volume conjecture which was first introduced by Rinat Kashaev.
We list here main references and observation on them.

\par\space{1em}
I first consider papers of Kashaev to recall his motivation on the Volume conjecture.

However, I misread \cite{Kashaev-1994-Qdilog6j} which does not deal with the Volume conjecture, I change the direction to understand the relation between $6j$-symbols and the combinatorics of 3-simplices.
In other words, I wanted to understand why the solution of the pentagon relation gives an assignment on labeled tetrahedra which is invariant under 2-3 Pachner moves, and that induces topological invariant (of glued 3-simplices).

The reference first mention about the relation between quantum dilogarith, $6j$-symbol and a Borel subalgebra of $U_q(\mathfrak{sl}_2)$, which also lives in my interest.
One of the goal of this paper is to state that the Quantized Roger's pentagon identity for the cyclic quantum dilogarithm $\Leftarrow$ coasociativity property of the Weyl algebra.
It seems that the author consider the solution of the pentagon identity via some module of the Weyl algebra (or its variation), and then consider state-sum invariant from coefficeints of this tensor.

\par\space{1em}

The basic motivation seems as follows.
The Clebsch--Gordan operator, whcih is a projection of simple modules of the Borel subalgebra of $U_q(\mathfrak{sl}_2)$, follows the simple rule.
\[
    V_p \otimes V_q \cong \bigoplus_{i=1}^N  V_{pq}
\]
where $V_p$ is a cyclic module where $E_N$ acts by the identity and $D_N$ acts by $z_p \in \mathbb{C}$.

By using this fact, one can assign the tensor $K_\alpha(p, q)$ to a trivalent graph whose edges are labeled by $p$, $q$, and $pq$ where $\alpha \in \mathbb{Z}/ N\mathbb{Z}$.
It induces an assignment on 3-valent graphs on $S^2$ (or maybe any graphs on the sphere more generally).
By dualizing graphs, we consider such graphs are boundary of polyhedra with triangular faces.
Hence, the assignment induces an tensor coefficient assignment on polyhedra with colored edges and colored vertices.
It seems that some properties of Clebsch--Gordan operators give the invariance of this assignment under gluing operation among polyhedra.
However, we ended the day during finding out which exectly property gives this invariance.


\bibliographystyle{amsalpha}
% \nocite{*}
\bibliography{bibliography}
\end{document}
